\documentclass[11pt]{article}
\usepackage[margin=1in]{geometry}
\usepackage{enumitem}
% =============================================================================
% Preambles/header.tex — shared LaTeX/Beamer preamble for this project
%
% Use from a Beamer lecture by prepending:
%     \input{header}
% and compiling with TEXINPUTS=../Preambles:$TEXINPUTS (the /compile-latex
% skill does this automatically).
%
% PALETTE CONTRACT. Color names defined here MUST match the names in
% Quarto/theme-template.scss so Beamer and Quarto renderings use the same
% palette. The scripts/check-palette-sync.sh script enforces this.
%
% To customize: edit the HEX values below AND the matching $variable in
% Quarto/theme-template.scss, then run scripts/check-palette-sync.sh.
% =============================================================================

% -----------------------------------------------------------------------------
% Engine detection + encoding
%
% Our /compile-latex skill uses XeLaTeX, where inputenc is unnecessary and
% can produce encoding warnings. Only load inputenc under pdfTeX.
% -----------------------------------------------------------------------------
\usepackage{iftex}
\ifPDFTeX
  \usepackage[utf8]{inputenc}
\fi

\usepackage{amsmath, amssymb, amsthm}
\usepackage{booktabs}
\usepackage{graphicx}

% -----------------------------------------------------------------------------
% PALETTE — mirrors Quarto/theme-template.scss variable names.
% Load xcolor and define colors BEFORE anything that references them
% (notably hyperref, hypersetup, and the Beamer color assignments).
% -----------------------------------------------------------------------------
\usepackage{xcolor}

\definecolor{primary-blue}{HTML}{012169}      % headings, accents
\definecolor{primary-gold}{HTML}{B9975B}      % emphasis, borders
\definecolor{highlight-yellow}{HTML}{F2A900}  % markers, alerts
\definecolor{light-bg}{HTML}{E8EDF5}          % title / section backgrounds
\definecolor{jet}{HTML}{1A1A1A}               % body text

% Semantic colors (used in diagrams and annotations)
\definecolor{positive}{HTML}{15803D}          % good / observed / identified
\definecolor{negative}{HTML}{B91C1C}          % bad / problematic / bias
\definecolor{neutral}{HTML}{525252}           % reference / context

% Highlight accents (match .hi-* classes in the SCSS)
\definecolor{hi-slate}{HTML}{314F4F}
\definecolor{hi-green}{HTML}{2E8B57}
\definecolor{hi-red}{HTML}{C41E3A}

% Semantic aliases so skill templates and snippets can write
% \textcolor{accent}{...} without hard-coding "primary-blue".
\colorlet{accent}{primary-blue}
\colorlet{accent2}{primary-gold}

% -----------------------------------------------------------------------------
% Hyperref — loaded AFTER xcolor + palette so link colors resolve.
% -----------------------------------------------------------------------------
\usepackage{hyperref}
\hypersetup{colorlinks=true, linkcolor=primary-blue, urlcolor=primary-blue,
            citecolor=primary-blue, pdfborder={0 0 0}}

% -----------------------------------------------------------------------------
% Course code — set once per deck (after \input{header}, before \begin{document})
% via \coursecode{CS401}. Shown in the footer on every slide so a deck's course
% is identifiable without opening the title page. Empty by default.
% -----------------------------------------------------------------------------
\makeatletter
\newcommand{\coursecode}[1]{\gdef\@coursecodetext{#1}}
\gdef\@coursecodetext{}
\makeatother

% -----------------------------------------------------------------------------
% Beamer theme assignments (only apply under Beamer).
%
% We detect Beamer via \@ifclassloaded{beamer}, which is the documented,
% non-internal way. This must be wrapped in \makeatletter/\makeatother
% because \@ifclassloaded contains an @-catcode character.
% -----------------------------------------------------------------------------
\makeatletter
\@ifclassloaded{beamer}{%
  \usefonttheme{professionalfonts}
  \setbeamercolor{structure}{fg=primary-blue}
  \setbeamercolor{frametitle}{fg=primary-blue}
  \setbeamercolor{title}{fg=primary-blue}
  \setbeamercolor{subtitle}{fg=primary-gold}
  \setbeamercolor{author}{fg=primary-blue}
  \setbeamercolor{institute}{fg=neutral}
  \setbeamercolor{date}{fg=primary-gold}
  \setbeamercolor{itemize item}{fg=primary-blue}
  \setbeamercolor{itemize subitem}{fg=primary-gold}
  \setbeamercolor{alerted text}{fg=primary-gold}
  \setbeamercolor{block title}{fg=white, bg=primary-blue}
  \setbeamercolor{block body}{bg=light-bg}
  % Semantic block variants: block=definition/concept (blue), exampleblock=
  % worked example (green/positive), alertblock=key takeaway or caution (gold).
  % Keep to <=2 colored boxes per slide across ALL three variants (INV-7).
  \setbeamercolor{block title example}{fg=white, bg=positive}
  \setbeamercolor{block body example}{bg=light-bg}
  \setbeamercolor{block title alerted}{fg=white, bg=primary-gold}
  \setbeamercolor{block body alerted}{bg=light-bg}

  % Minimal footer: course code (left) + page X / N (right), no navigation clutter
  \setbeamertemplate{navigation symbols}{}
  \setbeamertemplate{footline}{%
    \color{neutral}\footnotesize\hspace{0.7em}\@coursecodetext\hfill
    \insertframenumber/\inserttotalframenumber\hspace{0.7em}\vspace{0.3em}}
}{}
\makeatother

% -----------------------------------------------------------------------------
% TikZ setup — libraries the snippet gallery uses
% -----------------------------------------------------------------------------
\usepackage{tikz}
\usetikzlibrary{arrows.meta, positioning, calc, decorations.pathreplacing,
                fit, shapes.geometric, backgrounds}

% Shared TikZ styles — snippets and hand-written diagrams can pick these up
% via `\begin{tikzpicture}[every node/.style={...}]` or by naming specific
% styles (e.g., `\node[dag-node] (X) at (0,0) {X};`).
\tikzset{
  % Node styles for causal DAGs and similar
  dag-node/.style = {
    draw=primary-blue, fill=light-bg, rounded corners=2pt,
    minimum width=1.2cm, minimum height=0.9cm, align=center, font=\small
  },
  decision-node/.style = {
    draw=primary-gold, fill=white, diamond, aspect=1.8,
    minimum width=1.8cm, minimum height=1.2cm, align=center, font=\small,
    inner sep=1pt
  },
  % Edge styles — observed vs counterfactual
  observed-edge/.style = {-{Latex[length=2.5mm]}, thick, draw=jet},
  counterfactual-edge/.style = {-{Latex[length=2.5mm]}, thick, draw=neutral, dashed},
  confound-edge/.style = {-{Latex[length=2.5mm]}, thick, draw=negative, dashed},
  % Data-point styles for plots
  observed-dot/.style = {fill=primary-blue, draw=primary-blue, circle, inner sep=1.2pt},
  counterfactual-dot/.style = {fill=white, draw=neutral, circle, inner sep=1.2pt}
}

% -----------------------------------------------------------------------------
% Convenience macros
% -----------------------------------------------------------------------------
\newcommand{\muted}[1]{\textcolor{neutral}{#1}}
\newcommand{\key}[1]{\textcolor{primary-gold}{\textbf{#1}}}
\newcommand{\good}[1]{\textcolor{positive}{#1}}
\newcommand{\bad}[1]{\textcolor{negative}{#1}}

% Standout frame macro: full-bleed dark background for section transitions.
% Beamer-only (uses \begin{frame}/\setbeamercolor/\usebeamerfont) -- do not
% call from an article-class file (e.g. Notes/<CODE>/*-notes.tex). \muted,
% \key, \good, \bad, and \coursecode above are plain xcolor/gdef and are
% safe to use from either class; verified when Notes/article-class support
% was added.
% Expand: \transitionslide{Your transition title}
\newcommand{\transitionslide}[1]{%
  \begingroup
  \setbeamercolor{background canvas}{bg=primary-blue}
  \begin{frame}[plain]
    \centering
    \vfill
    {\usebeamerfont{title}\color{white}\Large #1\par}
    \vfill
  \end{frame}
  \endgroup
}

% Section divider frame (Beamer-only), an alternative to \transitionslide with a
% darker two-line style: near-black (jet) background, a small white label line
% above a larger gold title, and a thin gold rule along the bottom edge.
% Expand: \sectiondivider{Section 2}{Pointers}
\newcommand{\sectiondivider}[2]{%
  \begingroup
  \setbeamercolor{background canvas}{bg=jet}
  \begin{frame}[plain]
    \vspace*{3.0cm}
    {\color{white}\ttfamily\large #1\par}
    \vspace{0.5cm}
    {\color{primary-gold}\LARGE #2\par}
    \begin{tikzpicture}[remember picture, overlay]
      \draw[primary-gold, line width=2.5pt]
        ($(current page.south west)+(0,0.1)$) -- ($(current page.south east)+(0,0.1)$);
    \end{tikzpicture}
  \end{frame}
  \endgroup
}

% -----------------------------------------------------------------------------
% Channel watermark — "Concepts That Click" (YouTube).
%
% Article-class files (CompetitiveExam Q/A sets, Notes, Assignments) get a
% light diagonal draft watermark on every page plus a centered footer naming
% the channel. Beamer decks get a subtle TikZ background watermark instead
% (the footline already carries the course code). Centralized here so every
% MCQ PDF is branded without per-file edits.
% -----------------------------------------------------------------------------
\makeatletter
\@ifclassloaded{article}{%
  \usepackage{draftwatermark}
  \SetWatermarkText{Concepts That Click}
  \SetWatermarkScale{1}
  \SetWatermarkFontSize{2cm}
  \SetWatermarkLightness{0.86}
  \SetWatermarkAngle{45}
  \usepackage{fancyhdr}
  \pagestyle{fancy}
  \fancyhf{}
  \fancyfoot[C]{\footnotesize\textcolor{neutral}{Concepts That Click~|~YouTube: \href{https://www.youtube.com/@concepts.thatclick}{@concepts.thatclick}}}
  \renewcommand{\headrulewidth}{0pt}
  \renewcommand{\footrulewidth}{0pt}
  \fancypagestyle{plain}{%
    \fancyhf{}%
    \fancyfoot[C]{\footnotesize\textcolor{neutral}{Concepts That Click~|~YouTube: \href{https://www.youtube.com/@concepts.thatclick}{@concepts.thatclick}}}%
    \renewcommand{\headrulewidth}{0pt}%
    \renewcommand{\footrulewidth}{0pt}%
  }
}{}
\@ifclassloaded{beamer}{%
  \setbeamertemplate{background}{%
    \begin{tikzpicture}[remember picture, overlay]
      \node[rotate=45, opacity=0.055, align=center]
        at (current page.center)
        {\Huge\textcolor{neutral}{Concepts That Click}};
    \end{tikzpicture}%
  }
}{}
\makeatother 
\coursecode{JKSSB}
\renewcommand{\thesection}{32.\arabic{section}}
\newtheorem{example}{Example}[section]
\newenvironment{definitionbox}[1]{%
  \par\noindent\textbf{Definition (#1).}\ \itshape
}{\par\vspace{0.3em}}
\title{PowerPoint Animation and Transition --- Detailed Lecture Notes}
\author{JKSSB Computer Awareness}
\date{\today}

\begin{document}
\maketitle

\section{Animation and transition: the core distinction}
PowerPoint can make a presentation move, but it uses two different tools for
two different kinds of movement. An \textbf{animation} is an effect applied to
one object or to text \emph{inside} a single slide. A \textbf{transition} is an
effect applied to the movement \emph{from one slide to the next}. The simplest
way to keep them apart is to ask where the movement happens: movement inside a
slide is animation; movement between slides is a transition.

Both kinds of effect are added from the PowerPoint ribbon. Animation is
managed on the \textbf{Animations} tab, and transitions are managed on the
\textbf{Transitions} tab. Neither tool changes the content of the slides; they
change only how that content appears and how the deck moves from slide to
slide.

\begin{definitionbox}{Animation}
An effect applied to a single object or to text on a slide --- for example
making a picture appear, or making a title grow. Animation acts within a
slide.
\end{definitionbox}

\begin{definitionbox}{Transition}
An effect applied to the change from one slide to the next, such as Fade,
Push, or Morph. A transition acts between slides.
\end{definitionbox}

\begin{center}
\begin{tabular}{p{0.22\textwidth}p{0.32\textwidth}p{0.34\textwidth}}
\toprule
\textbf{Point} & \textbf{Animation} & \textbf{Transition} \\
\midrule
Where it applies & An object or text on a slide & The change between two slides \\
Ribbon tab & Animations & Transitions \\
Nearest confusion & Transitions & Entrance/Exit effects \\
\bottomrule
\end{tabular}
\end{center}

\section{The four animation effects}
When you select an object and open the Animations tab, PowerPoint groups the
effects into four categories. Each category answers a different question about
the object.

\begin{center}
\begin{tabular}{p{0.20\textwidth}p{0.64\textwidth}}
\toprule
\textbf{Effect} & \textbf{What it does} \\
\midrule
Entrance & Makes an object \textbf{appear} on the slide. \\
Exit & Makes an object \textbf{leave} the slide. \\
Emphasis & Draws attention to an object that is \textbf{already visible},
staying on the slide. \\
Motion Path & Moves an object \textbf{along a defined path}, such as a line or
curve. \\
\bottomrule
\end{tabular}
\end{center}

\textbf{Entrance} effects are used to reveal content step by step: the title or
a bullet point appears when the presenter is ready. \textbf{Exit} effects are
the opposite --- they take an object off the slide once it is no longer needed.
An entrance and an exit are not the same as an emphasis: an emphasis effect
does not add or remove the object, it only makes the object stand out while it
remains on the slide. A \textbf{motion path} is different again: it changes the
object's position, moving it from one place to another along a drawn path.

\begin{example}
A presenter wants a logo to travel from the left edge to the corner of the
slide. That is a \textbf{Motion Path}, because the object's position changes
during the show. If instead the logo simply spins once and stops in the same
place, that is an \textbf{Emphasis} effect. If the logo is not on the slide at
all until clicked, that is an \textbf{Entrance} effect.
\end{example}

Animations can be set to start when the slide appears or when the presenter
clicks, and each effect has a duration. The order and timing of the effects on
one slide are controlled from the \textbf{Animation Pane}.

\section{Transitions and Morph}
A transition plays as the presentation moves from the current slide to the
next. Common transitions include \textbf{Fade}, \textbf{Push}, \textbf{Wipe},
and \textbf{Dissolve}. Each transition can be given a duration and an optional
sound. From the Transitions tab, the \textbf{Apply To All} button copies the
chosen transition to every slide in the deck.

Transitions are often grouped by the strength of their effect. \textbf{Subtle}
transitions, such as Fade and Wipe, are gentle and professional.
\textbf{Exciting} transitions, such as Push and Split, use stronger movement.
\textbf{Dynamic Content} transitions respond to the slide content itself.

\begin{definitionbox}{Morph}
A modern transition that smoothly moves and animates objects as one slide
changes into the next. Morph works best when two slides share similar objects,
so the objects appear to glide from their old position to their new one.
\end{definitionbox}

\begin{example}
Slide 1 shows a circle on the left side; slide 2 shows the same circle on the
right side. With \textbf{Morph} applied to slide 2, the circle appears to
slide smoothly across the screen instead of jumping. Because Morph is a
transition, it is applied from the \textbf{Transitions} tab, not the Animations
tab.
\end{example}

\section{Starting and controlling the slide show}
The slide show has a small set of shortcut keys that JKSSB papers ask about
directly. The most important is \textbf{F5}, which starts the slide show from
the beginning --- that is, from the first slide. \textbf{Shift+F5} starts the
slide show from the current slide --- the one selected. \textbf{Alt+F5} starts
the presentation in \textbf{Presenter View}. \textbf{Esc} ends a running slide
show.

\begin{center}
\begin{tabular}{p{0.22\textwidth}p{0.60\textwidth}}
\toprule
\textbf{Key} & \textbf{Action} \\
\midrule
F5 & Start the slide show from the \textbf{beginning}. \\
Shift+F5 & Start the slide show from the \textbf{current} slide. \\
Alt+F5 & Start the presentation in \textbf{Presenter View}. \\
Esc & \textbf{End} the slide show. \\
\bottomrule
\end{tabular}
\end{center}

Once the show is running, more controls are available. The keys \textbf{N},
\textbf{Enter}, \textbf{Space}, or the right arrow advance to the next slide or
animation; \textbf{P} or the left arrow go back to the previous slide.
Pressing \textbf{B} turns the screen black and \textbf{W} turns it white, which
is useful for pausing discussion. Typing a slide number and pressing
\textbf{Enter} jumps straight to that slide.

\begin{example}
An item asks, ``Which key starts the slide show from the beginning?'' Trace
the key set: F5 is the full show from the first slide, Shift+F5 is from the
current slide, Alt+F5 opens Presenter View, and F7 is the spelling and grammar
check. The answer is \textbf{F5}. A distractor that offers F7 is testing
whether you confuse the show shortcut with spell check.
\end{example}

\section{Speaker notes, views and handouts}
The \textbf{Notes pane} sits below the slide in Normal view and holds
\textbf{speaker notes} --- the reminders the presenter reads while presenting.
Notes are not shown on the projected slide during a normal slide show; in
\textbf{Presenter View} the notes are visible only to the presenter, along with
the current slide, the next slide, and a timer.

PowerPoint provides several views for different jobs. \textbf{Normal view} is
used to edit one slide, with an outline on the side and the notes pane below.
\textbf{Outline view} works with the text of all slides as an outline.
\textbf{Slide Sorter view} shows small \textbf{thumbnails} of every slide and
is the view for rearranging slides. \textbf{Notes Page view} shows a slide
together with its speaker notes. \textbf{Reading View} reviews the show in a
window without going full screen, and \textbf{Slide Show view} presents the
deck full screen.

A \textbf{handout} is the printed version of the presentation that is given to
the \textbf{audience}. Handouts usually place several slides on one page ---
two, three, four, six, or nine per page --- and some handout layouts include
lined space for the audience to write. The key contrast is simple: the notes
pane is for the presenter, and a handout is for the audience.

\begin{definitionbox}{Handout}
A printed or exported layout of the presentation for the audience, normally
showing several slides on a single page, often with room for notes.
\end{definitionbox}

\section{SmartArt conversion and its limits}
\textbf{SmartArt} is a graphic that arranges text into a designed layout such
as a list, process, or hierarchy. A SmartArt graphic can be converted into
\textbf{shapes} or into \textbf{text}. Converting to \textbf{text} turns the
graphic into a plain bulleted list, so the graphical layout is lost. Converting
to \textbf{shapes} keeps individual shapes but removes the SmartArt
arrangement, and once the conversion is done the graphic cannot be turned back
into SmartArt --- the content that depended on the layout must be
\textbf{re-entered} or re-applied by hand.

\begin{example}
A student has a SmartArt process diagram and converts it to text to reuse the
wording in a document. The words survive, but the arrows, boxes, and arranged
layout do not. To get the diagram back, the SmartArt must be created again,
because a text conversion has dropped the design. The JA 2026 paper tested this
kind of limit --- which content cannot convert cleanly without re-entry --- as
its SmartArt item (PYQ-17).
\end{example}

\section{Exam retrieval and common confusions}
Five distinctions prevent most errors on this topic.
\begin{enumerate}
  \item Animation applies to an \textbf{object inside a slide}; a transition
    applies between \textbf{slides}. Movement inside a slide is animation;
    movement between slides is a transition.
  \item \textbf{Entrance}, \textbf{Exit}, \textbf{Emphasis}, and
    \textbf{Motion Path} are animation effects, not transitions.
  \item \textbf{F5} starts from the beginning; \textbf{Shift+F5} starts from
    the current slide; \textbf{Alt+F5} opens Presenter View; \textbf{Esc} ends
    the show. Do not confuse F5 with F7 (spell check) (TRAP-32).
  \item The \textbf{Notes pane} is for the presenter; a \textbf{handout} is for
    the audience; \textbf{Slide Sorter} is for reordering slides.
  \item \textbf{SmartArt} converted to text loses its graphical layout and may
    need the content re-entered.
\end{enumerate}

\begin{example}
An item lists four effects and asks which one is a transition. Read each
option against the definition: Entrance, Exit, and Emphasis all act on an
object within a slide, while Fade acts on the move between slides. Fade is the
transition. The same logic applies when an option offers Morph: Morph is a
transition, so it belongs with Fade, not with the animation effects.
\end{example}

\subsection*{Exercises}
\begin{enumerate}[label=\arabic*.]
  \item Distinguish an animation from a transition, giving one example of each.
  \item Name the four animation effect categories and state what each does.
  \item State the action of F5, Shift+F5, Alt+F5, and Esc during a slide show.
  \item Explain what the Notes pane is for and how a handout differs from it.
  \item What happens to a SmartArt graphic when it is converted to text?
\end{enumerate}

\subsection*{Solutions}
\begingroup\small
\begin{enumerate}[label=\arabic*.]
  \item An animation is an effect on one object or text inside a slide (for
    example, a picture entering). A transition is an effect on the change
    between two slides (for example, a Fade).
  \item Entrance makes an object appear; Exit makes an object leave; Emphasis
    draws attention to an object that stays on the slide; Motion Path moves an
    object along a defined path.
  \item F5 starts the show from the beginning; Shift+F5 starts from the current
    slide; Alt+F5 starts Presenter View; Esc ends the show.
  \item The Notes pane holds speaker notes for the presenter and is not shown
    to the audience during a normal show. A handout is the printed layout given
    to the audience, usually with several slides per page and room to write.
  \item It becomes a plain bulleted list; the graphical layout is lost, and the
    design must be re-applied or the SmartArt re-created.
\end{enumerate}
\endgroup
\end{document}
